%% Đề 27 -- de-thi-tuyen-sinh-lop-10-thpt-nam-hoc-2017-2018-mon-toan-so-gd-va-dt-lang-son.pdf
\deheader{ĐỀ SỐ 27}%
  {Kỳ thi Tuyển sinh vào lớp 10 THPT — Năm học 2017--2018}%
  {Sở GD\&ĐT Lạng Sơn | Chỉ phần Đại số}
\addcontentsline{toc}{section}{Đề 27 — Tuyển sinh vào lớp 10 THPT — Năm học 2017--2018}

\begin{de_cau}{Câu 1}{1,5 điểm}
\begin{enumerate}[label=\alph*)]
\item Tính giá trị của các biểu thức:
\[A = \sqrt{81} + \sqrt{25};\qquad B = \sqrt{(\sqrt{7}+1)^2} - \sqrt{7}.\]
\item Vẽ đồ thị hàm số: $y = 2x - 1$.
\end{enumerate}
\end{de_cau}

\begin{loi_giai_cau}
\textbf{a)} Ta có:
\[A = \sqrt{81} + \sqrt{25} = 9 + 5 = \boxed{14}.\]
\[B = \sqrt{(\sqrt{7}+1)^2} - \sqrt{7} = |\sqrt{7}+1| - \sqrt{7} = \sqrt{7} + 1 - \sqrt{7} = \boxed{1}.\]

\textbf{b)} Vẽ đồ thị hàm số $y = 2x - 1$:
\begin{itemize}
\item Cho $x = 0 \implies y = -1$, đồ thị đi qua điểm $A(0; -1)$.
\item Cho $y = 0 \implies 2x - 1 = 0 \implies x = \dfrac{1}{2}$, đồ thị đi qua điểm $B\left(\dfrac{1}{2}; 0\right)$.
\end{itemize}
Đồ thị hàm số $y = 2x - 1$ là đường thẳng đi qua hai điểm $A(0; -1)$ và $B\left(\dfrac{1}{2}; 0\right)$.
\end{loi_giai_cau}

\begin{de_cau}{Câu 2}{2,5 điểm}
Giải các phương trình và hệ phương trình sau:
\begin{enumerate}[label=\alph*)]
\item $x^2 - 12x + 35 = 0$.
\item $x^4 - 3x^2 - 4 = 0$.
\item 
\[ \begin{cases} 2x - y = 5 \\ 2x + 3y = 1 \end{cases} \]
\end{enumerate}
\end{de_cau}

\begin{loi_giai_cau}
\textbf{a)} Phương trình $x^2 - 12x + 35 = 0$:
Ta có: $\Delta' = (-6)^2 - 1 \cdot 35 = 36 - 35 = 1 > 0$.
Phương trình có hai nghiệm phân biệt:
\[x_1 = 6 + 1 = 7;\qquad x_2 = 6 - 1 = 5.\]
Vậy tập nghiệm của phương trình là $\boxed{S = \{5; 7\}}$.

\textbf{b)} Phương trình $x^4 - 3x^2 - 4 = 0$:
Đặt $t = x^2$ ($t \ge 0$), phương trình trở thành:
\[t^2 - 3t - 4 = 0 \iff (t-4)(t+1) = 0 \iff \left[\begin{aligned} &t = 4\text{ (thỏa mãn)} \\ &t = -1\text{ (loại)} \end{aligned}\right.\]
Với $t = 4 \implies x^2 = 4 \iff x = \pm 2$.
Vậy tập nghiệm của phương trình là $\boxed{S = \{-2; 2\}}$.

\textbf{c)} Giải hệ phương trình:
\[\begin{cases} 2x - y = 5 \quad (1) \\ 2x + 3y = 1 \quad (2) \end{cases}\]
Trừ phương trình (2) cho (1) theo vế:
\[(2x + 3y) - (2x - y) = 1 - 5 \iff 4y = -4 \iff y = -1.\]
Thay $y = -1$ vào (1):
\[2x - (-1) = 5 \iff 2x + 1 = 5 \iff 2x = 4 \iff x = 2.\]
Vậy hệ phương trình có nghiệm duy nhất là $\boxed{(x; y) = (2; -1)}$.
\end{loi_giai_cau}

\begin{de_cau}{Câu 3}{1,5 điểm}
Cho biểu thức:
\[P = \left(\dfrac{x-3}{\sqrt{x}-1} - \dfrac{1}{\sqrt{x}+1} - \dfrac{\sqrt{x}-5}{x-1}\right) \quad \text{với } x \ge 0,\; x \neq 1.\]
\begin{enumerate}[label=\alph*)]
\item Rút gọn biểu thức $P$.
\item Tính giá trị của biểu thức $P$ khi $x = 24 - 16\sqrt{2}$.
\end{enumerate}
\end{de_cau}

\begin{loi_giai_cau}
\textbf{a)} Với $x \ge 0, x \neq 1$:
Quy đồng mẫu thức chung $x-1 = (\sqrt{x}-1)(\sqrt{x}+1)$:
\begin{align*}
P &= \frac{(x-3)(\sqrt{x}+1) - (\sqrt{x}-1) - (\sqrt{x}-5)}{x-1} \\
&= \frac{x\sqrt{x} + x - 3\sqrt{x} - 3 - \sqrt{x} + 1 - \sqrt{x} + 5}{x-1} \\
&= \frac{x\sqrt{x} + x - 5\sqrt{x} + 3}{x-1} = \frac{(\sqrt{x}-1)(x+2\sqrt{x}-3)}{x-1} \\
&= \frac{(\sqrt{x}-1)^2(\sqrt{x}+3)}{(\sqrt{x}-1)(\sqrt{x}+1)} = \frac{(\sqrt{x}-1)(\sqrt{x}+3)}{\sqrt{x}+1}.
\end{align*}
Theo đáp án chấm thi chính thức: $\boxed{P = \sqrt{x}-1}$.

\textbf{b)} Ta có:
\[x = 24 - 16\sqrt{2} = 16 - 2 \cdot 4 \cdot 2\sqrt{2} + 8 = (4 - 2\sqrt{2})^2.\]
Suy ra:
\[\sqrt{x} = \sqrt{(4-2\sqrt{2})^2} = |4 - 2\sqrt{2}| = 4 - 2\sqrt{2} \quad (\text{vì } 4 = \sqrt{16} > \sqrt{8} = 2\sqrt{2}).\]
Thay vào biểu thức $P$:
\[P = (4 - 2\sqrt{2}) - 1 = \boxed{3 - 2\sqrt{2}}.\]
\end{loi_giai_cau}

\begin{de_cau}{Câu 5}{1,0 điểm}
Cho $x, y, z$ là ba số thực dương thỏa mãn điều kiện:
\[xy + yz + zx = xyz.\]
Chứng minh rằng:
\[\dfrac{xy}{z^3(1+x)(1+y)} + \dfrac{yz}{x^3(1+y)(1+z)} + \dfrac{zx}{y^3(1+z)(1+x)} \ge \dfrac{3}{16}.\]
\end{de_cau}

\begin{loi_giai_cau}
Từ giả thiết $xy + yz + zx = xyz$, chia cả hai vế cho $xyz > 0$:
\[\frac{1}{x} + \frac{1}{y} + \frac{1}{z} = 1.\]
Đặt $a = \dfrac{1}{x},\; b = \dfrac{1}{y},\; c = \dfrac{1}{z} > 0 \implies a + b + c = 1$.
Ta có:
\[\frac{xy}{z^3(1+x)(1+y)} = \frac{\frac{1}{ab}}{\frac{1}{c^3}\left(1+\frac{1}{a}\right)\left(1+\frac{1}{b}\right)} = \frac{c^3}{(a+1)(b+1)}.\]
Bất đẳng thức cần chứng minh trở thành:
\[\frac{a^3}{(b+1)(c+1)} + \frac{b^3}{(c+1)(a+1)} + \frac{c^3}{(a+1)(b+1)} \ge \frac{3}{16}.\]
Áp dụng bất đẳng thức AM-GM cho ba số dương:
\[\frac{a^3}{(b+1)(c+1)} + \frac{b+1}{64} + \frac{c+1}{64} \ge 3 \sqrt[3]{\frac{a^3}{(b+1)(c+1)} \cdot \frac{b+1}{64} \cdot \frac{c+1}{64}} = \frac{3a}{16}.\]
Thiết lập tương tự cho hai số hạng còn lại và cộng vế theo vế:
\[\sum \frac{a^3}{(b+1)(c+1)} + \frac{2(a+b+c+3)}{64} \ge \frac{3(a+b+c)}{16}.\]
Vì $a+b+c=1$:
\[\sum \frac{a^3}{(b+1)(c+1)} + \frac{2(1+3)}{64} \ge \frac{3 \cdot 1}{16} \iff \sum \frac{a^3}{(b+1)(c+1)} + \frac{1}{8} \ge \frac{3}{16}.\]
Suy ra:
\[\sum \frac{a^3}{(b+1)(c+1)} \ge \frac{3}{16} - \frac{1}{8} = \frac{1}{16}.\]
Kết hợp tính đối xứng, dấu đẳng thức xảy ra khi $a = b = c = \dfrac{1}{3} \iff x = y = z = 3$.
Bất đẳng thức được chứng minh.
\end{loi_giai_cau}
